\documentclass[12pt]{article}
\usepackage[spanish]{babel}
\usepackage[utf8]{inputenc}
\usepackage[T1]{fontenc}
\usepackage{amsmath, amssymb}
\usepackage{geometry}
\usepackage{enumitem}
\usepackage{needspace}
\usepackage{tikz}
\usetikzlibrary{babel}

\geometry{letterpaper, margin=2cm}
\setlength{\parindent}{0pt}
\setlength{\parskip}{6pt}

\newcommand{\puntografico}{0.35}
\newcommand{\opcionalgebra}[2]{\item[#1)] $#2$}

\begin{document}

\begin{center}
  {\Large \textbf{Borrador de items}}\\
  {\large Bloque 1. Geometria - Afirmacion 3}\\[4pt]
  Prueba Nacional Estandarizada de Matematica, Secundaria 2026
\end{center}

\section*{Afirmacion}
Resuelve problemas que involucran relaciones de posicion relativa entre rectas, rectas y circunferencias o puntos y circunferencias.

\section*{Items propuestos}

\begin{enumerate}[label=\textbf{\arabic*.}, itemsep=1.05cm]

\Needspace{0.92\textheight}
\item
En las instalaciones de un centro recreativo hay una fuente circular de $16$ m de diametro. Si se crea un sendero rectilineo que pasa por los puntos $(20,4)$ y $(20,14)$.

La siguiente representacion grafica, cuyas unidades estan en metros, muestra la ubicacion $O$ de la oficina y $C$ del centro de la fuente:

\begin{center}
\begin{tikzpicture}[scale=0.34]
  \path[use as bounding box] (-2,-1.5) rectangle (22,19);
  \draw[<->, thick] (-1,0) -- (21,0);
  \draw[<->, thick] (0,-1) -- (0,18);
  \node[anchor=west] at (21.3,-0.25) {$x$ (este)};
  \node[anchor=south east] at (1.1,18.3) {$y$ (norte)};
  \fill (0,0) circle (\puntografico) node[below left] {$O$};
  \draw[dashed] (10,0) -- (10,9);
  \draw[dashed] (0,9) -- (10,9);
  \draw[thick] (10,9) circle (8);
  \fill (10,9) circle (\puntografico) node[anchor=south east, inner sep=2pt] {$C$};
  \node[anchor=north] at (10,0) {$10$};
  \node[left] at (0,9) {$9$};
\end{tikzpicture}
\end{center}

De acuerdo con la informacion anterior, la recta que representa ese sendero, con respecto a la circunferencia que representa el borde de esa fuente, es

\begin{enumerate}[label=\Alph*)]
  \item exterior.
  \item secante.
  \item tangente.
  \item Ninguna de las anteriores.
\end{enumerate}

\Needspace{0.96\textheight}
\item
Un juego consiste en lanzar una ficha desde cierta distancia hacia un objetivo circular colocado en el piso de un salon de actividades. La medida del radio de ese objetivo es $5$ m.

La siguiente representacion grafica, cuyas unidades estan en metros, muestra la circunferencia que representa el borde del objetivo:

\begin{center}
\begin{tikzpicture}[scale=0.38]
  \path[use as bounding box] (-2,-1.5) rectangle (17,16);
  \draw[<->, thick] (-1,0) -- (16,0);
  \draw[<->, thick] (0,-1) -- (0,15);
  \node[anchor=west] at (16.3,-0.25) {$x$ (este)};
  \node[anchor=south east] at (1.1,15.3) {$y$ (norte)};
  \draw[dashed] (9,0) -- (9,13);
  \draw[dashed] (0,8) -- (9,8);
  \draw[dashed] (0,13) -- (9,13);
  \draw[thick] (9,8) circle (5);
  \fill (9,8) circle (\puntografico);
  \fill (9,13) circle (\puntografico);
  \node[anchor=north] at (9,0) {$9$};
  \node[left] at (0,8) {$8$};
  \node[left] at (0,13) {$13$};
\end{tikzpicture}
\end{center}

Ademas, durante el juego:

\begin{itemize}
  \item Mariela lanzo una ficha y esta se ubico en el punto correspondiente a $(12,12)$.
  \item Andres lanzo una ficha y esta se ubico en el punto correspondiente a $(4,8)$.
  \item Valeria lanzo una ficha y esta se ubico en el punto correspondiente a $(15,9)$.
  \item Esteban lanzo una ficha y esta se ubico en el punto correspondiente a $(8,4)$.
\end{itemize}

De acuerdo con la informacion anterior, \textquestiondown cual de esas personas lanzo la ficha que se ubico en el exterior de ese objetivo?

\begin{enumerate}[label=\Alph*)]
  \item Mariela.
  \item Andres.
  \item Valeria.
  \item Esteban.
\end{enumerate}

\Needspace{0.96\textheight}
\item
Una rotonda tiene forma circular y la medida de su diametro es $24$ m. Las rectas $h$ y $j$ representan senderos que pasan por el centro de la rotonda. Las rectas $k$ y $n$ representan otros senderos del lugar. La siguiente representacion grafica, en la que las unidades estan en metros, muestra la rotonda y esos senderos:

\begin{center}
\begin{tikzpicture}[scale=0.27]
  \path[use as bounding box] (-10,-7) rectangle (38,33);
  \draw[<->, thick] (-8,0) -- (36,0);
  \draw[<->, thick] (0,-5) -- (0,31);
  \node[anchor=west] at (36.7,-0.5) {$x$};
  \node[anchor=south east] at (1.6,31.7) {$y$};
  \draw[thick] (18,14) circle (12);
  \draw[very thick, <->] (18,-4) -- (18,32);
  \draw[very thick, <->] (2,14) -- (34,14);
  \draw[very thick, <->] (1,4) -- (35,4);
  \draw[very thick, <->] (7,29) -- (31,-2);
  \draw[thick] (18,4) -- (16.4,4) -- (16.4,5.6) -- (18,5.6);
  \draw[dashed] (0,14) -- (18,14);
  \draw[dashed] (18,0) -- (18,14);
  \fill (18,14) circle (\puntografico);
  \node[anchor=north] at (18,0) {$18$};
  \node[left] at (0,14) {$14$};
  \node[left] at (0,4) {$4$};
  \node[anchor=south east] at (18,32) {$h$};
  \node[anchor=south west] at (34,14) {$j$};
  \node[anchor=south east] at (8,29) {$k$};
  \node[anchor=south west] at (34,4) {$n$};
\end{tikzpicture}
\end{center}

De acuerdo con la informacion anterior, \textquestiondown cuales de esas rectas, que representan senderos, son perpendiculares entre si?

\begin{enumerate}[label=\Alph*)]
  \item $j$ y $k$.
  \item $h$ y $j$.
  \item $k$ y $n$.
  \item $j$ y $n$.
\end{enumerate}

\Needspace{0.92\textheight}
\item
Una antena de telecomunicaciones emite una senal que permite a los telefonos celulares, que se encuentran a $6$ km o menos alrededor de esa antena, realizar y recibir llamadas. A continuacion, se muestra la representacion grafica, en la que las unidades estan en kilometros, de la circunferencia que corresponde al alcance maximo de la senal que emite esa antena:

\begin{center}
\begin{tikzpicture}[scale=0.38]
  \path[use as bounding box] (-2,-1.5) rectangle (17,16);
  \draw[<->, thick] (-1,0) -- (16,0);
  \draw[<->, thick] (0,-1) -- (0,15);
  \node[anchor=west] at (16.3,-0.25) {$x$};
  \node[anchor=south east] at (1.1,15.3) {$y$};
  \draw[thick] (7,8) circle (6);
  \draw[dashed] (7,0) -- (7,14);
  \draw[dashed] (0,8) -- (13,8);
  \draw[dashed] (0,14) -- (7,14);
  \draw[dashed] (13,0) -- (13,8);
  \fill (7,8) circle (\puntografico);
  \fill (7,14) circle (\puntografico);
  \fill (13,8) circle (\puntografico);
  \node[anchor=north] at (7,0) {$7$};
  \node[anchor=north] at (13,0) {$13$};
  \node[left] at (0,8) {$8$};
  \node[left] at (0,14) {$14$};
\end{tikzpicture}
\end{center}

De acuerdo con la informacion anterior, si las ubicaciones que tienen los telefonos celulares $W$, $K$, $L$ y $M$ corresponden, respectivamente, a los puntos $(10,12)$, $(14,8)$, $(2,2)$ y $(7,15)$, entonces \textquestiondown en cual de esos telefonos se pueden realizar y recibir llamadas por medio de la senal de esa antena?

\begin{enumerate}[label=\Alph*)]
  \item $W$
  \item $K$
  \item $L$
  \item $M$
\end{enumerate}

\end{enumerate}

\section*{Clave y justificacion breve}

\begin{enumerate}[label=\textbf{\arabic*.}]
  \item \textbf{A.} El sendero esta representado por la recta $x=20$. La distancia del centro $(10,9)$ a esa recta es $10$, mayor que el radio de la circunferencia, que mide $8$ m; por eso la recta es exterior.
  \item \textbf{C.} El centro es $(9,8)$ y el radio es $5$. Para Valeria, $(15-9)^2+(9-8)^2=37>25$, por lo que su ficha se ubico en el exterior. Las demas fichas estan en el interior o sobre la circunferencia.
  \item \textbf{B.} La recta $h$ es vertical y la recta $j$ es horizontal; por lo tanto, se intersecan formando angulos rectos y son perpendiculares.
  \item \textbf{A.} El centro de la circunferencia es $(7,8)$ y el radio es $6$. Para $W(10,12)$ se cumple $(10-7)^2+(12-8)^2=25\leq 36$, por lo que esta dentro del alcance. Los telefonos $K$, $L$ y $M$ quedan a una distancia mayor que $6$ km.
\end{enumerate}

\end{document}
